
faisant intervenir les anneaux de classes de cycles (mod. équivalence rationnelle) $A(X)$ et $A(Y)$, commutativité établie dans SGA 5 IV.

\textbf{6.3.} Notons que dans chacun des cas envisagés plus haut où on arrive à définir un homomorphisme de Gysin, on obtient en fait une construction pour la cohomologie à coefficients dans $\Zell(*)$, et non seulement dans $\Qell(*)$. Cela permet, lorsque $f$ est une \uline{immersion}, de poser la question de la validité de la formule de Riemann--Roch ``sans dénominateurs'', comme il est expliqué au n° 3. Lorsque $X$ et $Y$ sont tous deux lisses sur un corps $k$, la commutativité de (6.4) nous montre qu'il suffirait de prouver la formule sans dénominateurs dans l'anneau de Chow $A(Y)$ cf. (3.3).

\textbf{6.4.} Lorsqu'on ne fait pas d'hypothèse de régularité sur $Y$, et même si $X$ et $Y$ sont de type fini sur le corps $\C$, on ne voit plus de définition plausible d'un homomorphisme de Gysin cohomologique (6.2), même pour une immersion régulière de codimension $1$. Peut-on en trouver un néanmoins avec des propriétés raisonnables, en particulier rendant commutatif le carré (6.3)? La question est liée évidemment à celle de définir, pour un $X$ d'intersection complète de codimension $i$ dans $Y$, ``sa classe de cohomologie''
\[
 \gamma(X)\in H^{2i}(Y,\Zell(i)),
\]
de façon à transformer intersections (sans composantes excédentaires) en cup-produits. Cette dernière question se résout affirmativement tout au moins modulo torsion, i.e. dans $H^{2*}(Y,\Qell(*))$, en définissant
\[
 \gamma(X)=1/(-1)^{i-1}(i-1)!\,c_i(\underline{O}_X).
\]
Mais j'ignore si $c_i(\underline{O}_X)$ est divisible par $(i-1)!$ dans $H^{2i}(Y,\Zell(i))$. En particulier, si $Y$ est un schéma projectif de dimension $3$ sur le corps des complexes, disons, et $X$ réduit à un point (singulier) de $Y$, $c_3(\underline{O}_X)$ est-il divisible par $2$ dans $H^6(Y,\Z_2(3))$, et en particulier la classe de Chern est-elle nulle dans $H^6(Y,\Z/2\Z)$? Signalons également dans la même voie la question suivante: soit $Y$ un schéma projectif (singulier) de dimension $4$, $X$ un sous-schéma fermé de dimension $1$ qui soit globalement une intersection de trois diviseurs positifs $D_1,D_2,D_3$, considérons le produit de leurs classes de Chern mod. $2$:
\[
 c_1(D_1)c_1(D_2)c_1(D_3)\in H^6(Y,\Z/2\Z);
\]
